\honest{Beware of Other-Optimizing}{Eliezer Yudkowsky}{2009}

I have noticed that aspiring rationalists ``vastly overestimate their ability to optimize other people's lives,'' and I think I know why.

You read nineteen webpages of self-improvement advice. Most of it seems silly, and the two methods you try fail. Then the twentieth works, ``STARS ABOVE,'' and now you know ``the real way.'' You know most people do not listen, but this one is a friend you trust to listen. So you take a friend by the shoulder and tell them how to do it.

People do this to me a lot. They know not to tell their families, but Eliezer Yudkowsky needs it and might understand. A board member of my institute told me I did not need a raise to keep up with inflation, because an online coupon service would cut my food bill. I believed it, because a trusted friend said it with confidence. My girlfriend tried the service and soon gave up. \nb{The coupon tip came from someone with a say in the author's pay, it seems, and was given as a reason against a raise. So the first example is already a case of the leverage the post warns about later.} Had I read the same tip on a blog, even one by Scott Aaronson, I would have moved on; delivered in person, it felt like being told the secret. It took me a while to see that. Personal advice is ``no more and no less likely to work for me'' than a blog post by any intelligent writer.

My advisers' methods are as varied as the productivity blogs. I take that to mean that people differ in what works for them; my advisers take it to mean that most advice is bad, except theirs. Some advice does help me, such as ``Stuck In The Middle With Bruce.'' \nb{A LessWrong post by CronoDAS from the day before, introducing an essay on Magic: the Gathering about ``people's need to lose.''}

``Different things work for different people'' gives me a squicky feeling, since Dark Side Epistemology uses the phrase to shield claims from criticism, much like ``Different things are true for different people,'' which is simply false. But until someone finds the general laws of what works for whom, we mess with surface tricks that work for one person and not another, and the best we can do is accept ``No'' for an answer.

This matters most if you have power over the other person. Power is easy to abuse without noticing, and the safeguards work only if you actually use them. I recall a post showing that power reduces empathy, though I cannot find it. \nb{Galinsky and colleagues (2006) found that power reduced perspective taking. Schmid Mast, Jonas and Hall, later in 2009, found in four studies that ``high power resulted in more interpersonal sensitivity than low power.''}

If your method needs willpower, you will call the other person lazy and push. Sometimes, I suppose, they are. But the ability to tell when more is within someone's capacity, without burning them out, is ``extremely rare.'' Bosses who have it are worth their weight in silver, and I do not have it.

Do not assume you have it because your intentions are good, because you would not do to others what you would not want done to you, or because no one has complained. Maybe they are scared. I have seen a rationalist, obviously powerful to me, amazed to learn he might be feared. Be careful whenever you hold leverage, such as a decision or something the other person needs. Ayn Rand's rule over her followers, I suggest, came from power and leverage she could not resist using to optimize others.

``We underestimate the distance between ourselves and others,'' not only in inferential distance but in temperament, ability, situation and resources, unspoken knowledge, unnoticed skills, luck and interior landscape. Even I am surprised when what worked for me fails for someone else, but having been optimized so often, I recognize the distance when it hits me. \nb{Research agrees. Ross, Greene and House named it the false consensus effect in 1977: people see their own choices as common and appropriate.}

Maybe being pushed works for you, and you do not feel sick when someone with power over you reorganizes your life. I do not know what makes you tick. All I can do is accept it when someone says my method does not work for them, and keep looking for the deeper laws. \nb{The next post in the book, ``Practical Advice Backed By Deep Theories,'' takes up that search.}
